\scanpage{46}
\mspage{37}

Je vais reprendre ici la prop.~1 p.~1 de CF \uline{V}, dans un contexte ``indiciel'' qui sera, en fait, un peu plus général.

\uline{Proposition} Soient $I$ un ens. ordonné, \deleted{\unsure{$P$ un}} $X$ un ensemble (appelé \unsure{$A$} dans loc. cit, et qui devient $P$ dans les pages précédentes), $\Phi=(X_i)_{i\in I}$ une famille d'ens indexée par $I$. \inserted{On suppose $i\leq j\Rightarrow X_i\subset X_j$.}
% margin note written slanted
\margin{[Je ne suis pas sûr si je l'utilise vraiment, si on ne le \unsure{suppose} pas, pour définir $X_i^{\circ}$ \unsure{(?)} il faut prendre $X_i\setminus X_i\cap\partial X_i$ (car on ne sait pas si $\partial X_i\subset X_i$)}
On pose, $\forall i\in I$
\[
\begin{aligned}
\partial X_i&=\bigcup_{j<i}X_i\subset X_i\\
X_i^{\circ}&=X_i\setminus\partial X_i\\
|\Phi|&=\bigcup_{i\in I}X_i.
\end{aligned}
\]
On note que $\forall x\in X$, \deleted{$i\in I$,} si on pose
\[
I^x=\{i\in I\mid x\in X_i\}
\]
on a l'équivalence
\[
x\in X_i^{\circ}\Longleftrightarrow i\text{ est un élément minimal de }I^x.
\]

Conditions équivalentes

(i) $\forall x\in|\unsure{\Phi'|\Phi}|$, l'ensemble $I^x$ a un plus petit élément.

(ii) Pour tout $x\in|\Phi|$, l'ensemble $I^x$ a un élément minimal, \inserted{(autom. si $I$ fini)}, et de plus pour $i,j\in I$, on a
\[
X_i\cap X_j=\bigcup_{\unsure{k|\underset{\sim}{k}}\leq i,j}X_k.
\]

(iii) On a les conditions suivantes

\begin{enumerate}[label=\alph*)]
\item $\forall(i,j)\in I$ avec $i\ne j$, on a
\[
X_i^{\circ}\cap X_j^{\circ}=\varnothing
\]
\item \deleted{$\unsure{\forall x\in X}$} $\forall i\in I$, on a
\[
X_i=\bigcup_{j\leq i}X_j^{\circ}
\]
\margin{(autom. si $I$ fini)}
\end{enumerate}

\deleted{\unsure{NB} a) et b) \unsure{impliquent} que la famille \unsure{sera} pour}
\inserted{(iv) Posant}
\[
I'=\{i\in I\mid X_i^{\circ}\ne\varnothing\},
\]
\deleted{alors la famille des} On veut :

$\alpha)$ La famille $\{X_i^{\circ}\}_{i\in I'}$ forme une \uline{partition} de $|\Phi|$, \deleted{\unsure{et}}
% archive sheet number: 46
